% ===========================================================================
%  sviluppatore/05-estendere/09-importare.tex
% ===========================================================================
\section{Importare un artefatto scritto altrove}\label{sez:importare}

Un artefatto scritto fuori dall'app non va copiato a mano nella cartella
giusta: in Studio ogni famiglia ha il pulsante d'importazione accanto alla
sua ricerca, nella colonna di sinistra (\vediSez{studio-banco}); accetta un
file \file{.py}, una cartella o uno \file{.zip} e ci cerca dentro la forma
della categoria. La logica è in \file{dashboard/studio\_import.py}, le rotte in
\file{dashboard/ui\_ext\_import.py}, la finestra in coda a
\file{static/js/studio.js}.

\begin{center}
\begin{tabular}{@{}L{24mm}L{30mm}L{84mm}@{}}
\toprule
\textbf{Categoria} & \textbf{Si sceglie} & \textbf{Che cosa diventa un artefatto} \\
\midrule
\rigaTenue
strategia, approccio & cartella, zip, .py
  & ogni cartella che contiene \file{strategy.py} / \file{approach.py}, con
    tutti i file e le sottocartelle (di due annidate vale la più esterna);
    un \file{.py} che assegna \cod{DESCRIPTOR} / \cod{APPROACH\_DESCRIPTOR},
    attorno al quale si crea la cartella \\
overlay, weight mode & cartella, zip, .py
  & la cartella con \file{overlay.py} / \file{weight\_mode.py}; un \file{.py}
    con \cod{META} (dizionario di letterali con \cod{key}) e \cod{def compute},
    che diventa la cartella con il nome della chiave \\
\rigaTenue
metrica & solo .py
  & ogni file è una metrica, con il nome del file ridotto alle regole
    (\cod{Live VS backtest.py} \textrightarrow{} \cod{live\_vs\_backtest});
    una cartella o uno zip si rifiutano con il perché \\
\bottomrule
\end{tabular}
\end{center}

\subsection{Due tempi}

\simbolo{analizza} legge la scelta senza eseguire niente. Ogni \file{.py} si
analizza con \cod{ast}: la sintassi (un errore ferma l'artefatto, con file e
riga), il descrittore o il \cod{META}, i campi che si leggono senza eseguire
--- \cod{name}, \cod{label}, \cod{description}, \cod{backtest\_script},
anche quando sono f-string di costanti del modulo
(\cod{f"strategies/\{NOME\}/run\_backtest.py"}). Gli artefatti buoni si
mettono da parte in \file{.studio\_import/<gettone>/} con un
\file{manifest.json}; la risposta è una scheda per artefatto, con stato
(\cod{new}, \cod{exists}, \cod{blocked}), errori, avvisi e note.
\simbolo{installa} sposta al loro posto quelli confermati, con
\cod{os.replace}: la cartella di lavoro sta sotto la radice dell'app, sullo
stesso disco, e lo spostamento è atomico.

\begin{description}
  \item[Il nome] è quello della cartella, perché il registro usa quello
    (\vediSez{due-forme}). Per strategie e approcci deve rispettare la
    regola dell'editor, \cod{\^{}[a-z][a-z0-9\_]\{0,49\}\$}: una cartella
    \file{My-Strategy} si ferma, perché rinominarla romperebbe il codice che
    costruisce percorsi dal proprio nome (\cod{backtest\_script}). Un nome
    diverso da quello del descrittore è un avviso. Con i file in cima alla
    scelta (uno zip senza cartella, più \file{.py} insieme) vale il nome dello
    zip, o quello del descrittore se lo zip ha un nome qualsiasi.
  \item[I file aggiunti] \file{\_\_init\_\_.py} se manca, e il
    \file{run\_backtest.py} che scrive l'editor quando il descrittore lo
    dichiara e non c'è.
  \item[Ciò che c'è già] un artefatto con lo stesso nome si sostituisce solo
    se scelto; il precedente va in
    \file{.studio\_import/sostituiti/<categoria>/<nome>-<ora>/} (le ultime
    cinque versioni). Una strategia e un approccio con lo stesso nome
    avrebbero gli stessi risultati: si ferma. Un'obsoleta con lo stesso nome
    (\file{\_nome}) e risultati rimasti sono note.
  \item[Chi lo usa] non si sostituisce ciò che un runner del Live usa o che
    un backtest sta leggendo: la stessa regola di «nascondi» ed «elimina»
    (\simbolo{strategie\_obsolete.uso}).
  \item[Dopo] i moduli dell'artefatto escono da \cod{sys.modules}, il
    registro della categoria lo ricarica (\simbolo{reload\_registry} per
    overlay e weight mode) e la risposta dice se lo carica, o il perché
    (\cod{ModuleNotFoundError: No module named 'sklearn'}). Il
    ripopolamento dell'app installata non lo tocca più
    (\simbolo{core.workspace.escludi}): un artefatto importato è dell'utente.
  \item[Una metrica] si prova come con \voce{Validate}
    (\simbolo{prova\_della\_metrica}): sullo stream con più dati, se l'app lo
    dà, altrimenti su un'equity finta. Un errore è un avviso, perché può
    dipendere dai dati (la metrica del confronto con il live si ferma se
    \cod{START} è dopo i dati della prova).
\end{description}

\subsection{Uno zip}

Si legge in memoria con dei limiti: \num{5000} file, \num{200} MB
decompressi, \num{100} MB per file, \num{120} MB di richiesta. Un percorso
assoluto, con \cod{..} o una lettera di unità, e un collegamento simbolico
fermano tutto lo zip. \file{\_\_pycache\_\_}, \file{.pyc}, \file{\_\_MACOSX},
\file{.git}, \file{.DS\_Store} e simili restano fuori, e il browser non li
manda nemmeno. Il resto entra com'è: un modello in \file{.pkl} o un
\file{README.md} sono parte dell'artefatto.

\begin{attenzione}{Il codice importato non è filtrato}
Il validatore che rifiuta file, rete e comandi vale per il codice che genera
l'agente e per i blocchi. Un artefatto importato è codice dell'utente, come
quello scritto nell'editor di Studio, che controlla solo la sintassi: la
porta dei dati di una strategia deve poter scaricare. Lo dice anche la guida
utente: si importa solo codice di cui ci si fida.
\end{attenzione}

\subsection{Le rotte}

\begin{center}
\begin{tabular}{@{}L{70mm}L{68mm}@{}}
\toprule
\textbf{Rotta} & \textbf{Che cosa fa} \\
\midrule
\rigaTenue
\cod{GET /studio/api/import/formats} & che cosa accetta ogni categoria e come
  lo spiega la finestra (\simbolo{formati}) \\
\cod{POST /studio/api/import/<categoria>} & multipart: \cod{kind}
  (\cod{py}, \cod{folder}, \cod{zip}), \cod{files}, \cod{paths} (i percorsi
  relativi di una cartella, nello stesso ordine), \cod{source}. Risponde con
  le schede e il gettone; con \cod{ok: false} se dentro non c'è niente della
  categoria \\
\rigaTenue
\cod{POST \ldots/<categoria>/confirm} & \cod{\{token, items: [\{id, replace\}]\}}:
  gli esiti, uno per artefatto \\
\cod{POST \ldots/<categoria>/cancel} & \cod{\{token\}}: la scelta messa da
  parte si butta (le altre si buttano dopo sei ore) \\
\bottomrule
\end{tabular}
\end{center}

\noindent
Con l'import, \file{studio\_browser.py} tratta overlay e weight mode come
forma «mista», cartella o file singolo: gli elenchi di overlay e weight mode
erano vuoti, perché contavano solo i file singoli mentre gli store scrivono la
cartella da tempo.
